\section{Discussion}
\label{sec:discussion}

\subsection{Trapdoor properties and confidentiality}

The BHF was originally motivated by \emph{trapdoor} data structures for
encrypted search~\cite{es}.
A cipher set (or map) is one whose binary representation reveals minimal
information about the underlying data beyond what is accessible through the
defined interface.

The maximum entropy property (\cref{thm:max_entropy}) is central to the
trapdoor property: the BHF encoding is incompressible (conditioned on its length),
so no side-channel information beyond the encoding length is available to an
adversary.
The encoding length itself reveals only an estimate of the cardinality, and
this leakage can be controlled via the entropy--space tradeoff
(\cref{sec:entropy}).

For Bernoulli maps, the choice of value encoder affects the trapdoor property.
An optimal (e.g., Huffman) encoder minimizes space but reveals distributional
information about the values.
A uniform-length encoder preserves the trapdoor property at the cost of higher $\mu$.

\subsection{Comparison with existing data structures}

\paragraph{Bloom filter~\cite{bloom1970,bf_survey}.}
The Bloom filter is a practical implementation of the positive Bernoulli set
($\fnrate = 0$).
It achieves $\log_2(e) \cdot \log_2(1/\fprate)$ bits per element, which
exceeds the information-theoretic lower bound of $-\log_2 \fprate$ by a
factor of $\log_2 e \approx 1.44$.
The Bloom filter does not support false negatives ($\fnrate > 0$) or maps.

\paragraph{Perfect Hash Filter~\cite{phf}.}
The PHF uses a perfect hash function to achieve $\log_2(e/\fprate)$ bits per
element for positive Bernoulli sets.
This is closer to optimal than the Bloom filter but still exceeds the lower
bound.
The Perfect Map Filter~\cite{pmf} extends this to maps.

\paragraph{Bernoulli Hash Function (equality/threshold).}
The BHF matches the lower bound ($-\log_2 \fprate + \mu$ bits per element) to
leading order in expectation, and its salt is uniform given its length.
However, the construction time is exponential in $m$ (expected
$\mathcal{O}(1/p)$ trials), so it is an illustration of a known bound, not a
practical structure; the bound was already known to be achievable with
efficient construction~\cite{pagh2005,dietzfelbinger2008retrieval}.

\paragraph{Bernoulli Hash Function (adaptive threshold).}
The adaptive variant (\cref{sec:adaptive}) has a random FPR
($\fprate \approx \betadist(p, m - p + 1)$) and needs no salt search for sets,
but its recall is about $p/m$ while its mean FPR is at least $p/(m+1)$
(\cref{sec:adaptive_contract}).
Its $\mathcal{O}(\log N / m)$ per-element cost reflects this weaker contract,
not the randomness of the FPR (\cref{rem:lb_not_violated}).
At a fixed recall and FPR cap it reduces to fixed-threshold search
(\cref{thm:price_certainty}).
It is not a useful membership structure.

\paragraph{Cuckoo filter~\cite{cuckoo_filter}.}
The cuckoo filter stores fingerprints in a cuckoo hash table, achieving
approximately $(2 + 3/\log_2(1/\fprate)) \log_2(1/\fprate)$ bits per element
with $\mathcal{O}(n)$ construction time and support for deletion.
It exceeds the lower bound by a small multiplicative factor.

\paragraph{XOR and binary fuse filters~\cite{xor_filter}.}
XOR filters achieve approximately $1.23 \log_2(1/\fprate)$ bits per element
with $\mathcal{O}(n)$ construction time, coming within 23\% of the lower bound.
Binary fuse filters further reduce this overhead.
Both are immutable (no deletion or insertion after construction).

\paragraph{Ribbon filter~\cite{ribbon_filter}.}
The ribbon filter uses a banded linear system to achieve approximately
$\log_2(1/\fprate) + 1.5$ bits per element with $\mathcal{O}(n)$ construction
time, approaching the lower bound with a small additive overhead.

\paragraph{Optimal construction and seed search.}
The space lower bound of $-\log_2 \fprate$ bits per element is due to Carter
et al.~\cite{carter1978} (with false negatives, Pagh and
Rodler~\cite{pagh2001lossy}).
Pagh, Pagh, and Rao~\cite{pagh2005} showed it is achievable up to lower-order
terms with efficient construction, and Dietzfelbinger and
Pagh~\cite{dietzfelbinger2008retrieval} gave succinct retrieval structures and
the general reduction from retrieval to approximate membership, which covers
maps.
The Bloomier filter~\cite{chazelle2004bloomier} already stores static
functions and accepts a query when the decoded value lies in a valid range.
Brute-force seed search near the information-theoretic scale is the
mechanism of minimal perfect hash functions such as
RecSplit~\cite{esposito2020recsplit} and ShockHash~\cite{lehmann2024shockhash};
Lehmann, Sanders, Walzer, and Ziegler~\cite{lehmann2025seeds} analyze one
seed that satisfies independent per-key constraints, which is the
zero-false-negative BHF, and improve on it by encoding sequences of seeds
jointly.
The BHF reaches the same bound by a simpler but exponential-time exhaustive
search.

\subsection{Practical considerations}

While the BHF construction is exponential-time and thus impractical for large
sets, it serves several purposes:
\begin{enumerate}
    \item It gives a short, self-contained illustration that the
    information-theoretic lower bound is achievable in expectation (a fact
    already established by the constructions above).
    \item It provides a benchmark against which practical data structures
    (Bloom filters, PHFs) can be compared.
    \item For very small sets ($m \leq 10$), the construction is feasible and
    may be useful in specialized applications such as encrypted search indices.
    \item The adaptive threshold variant runs in $\mathcal{O}(m \log m)$ time
    for any set size, but because its recall is about $p/m$ at mean FPR about
    $p/(m+1)$ it does not provide a useful membership test.
\end{enumerate}

\subsection{Encrypted search application}

In encrypted search~\cite{es}, a client constructs BHF instances of document
term sets and stores them on an untrusted server.
To search, the client provides the server with a query term and the server
tests membership in each BHF without learning the document contents.
The Bernoulli map variant enables rank-ordered search: each key (term) maps
to a relevance score, and the server can retrieve approximate scores for
query terms.
The maximum entropy and trapdoor properties ensure that the server learns
only what is revealed by the query results.
